\chapter{总结与展望}
\label{chap:conclusion}

\section{本文总结}

本文面向生成式 Next POI 推荐任务中的地理信息建模问题展开研究，重点关注现有基于 Semantic ID 的方法在离散表示学习阶段对空间关系刻画不足这一局限。随着大语言模型和生成式推荐方法的发展，基于离散语义标识的生成式 POI 推荐逐渐成为研究热点。然而，现有方法通常将地理信息更多作为预测阶段的外部辅助约束，较少在 Semantic ID 的形成过程中显式考虑 POI 之间的真实空间关系，导致离散表示虽然能够较好保留语义信息，却难以充分反映地理邻近性与区域结构特征。针对这一问题，本文在 GNPR-SID 框架基础上，提出了一种面向表示学习阶段的隐式地理信息注入方法，尝试从离散表示构建源头提升 POI 表示质量，从而改善生成式 Next POI 推荐效果。

实验部分基于 NYC、TKY 和 CA 三个公开数据集，对本文方法进行了系统验证。结果表明，与基线方法相比，本文方法在三个数据集上的 Acc@1 均取得了一定提升，其中在 NYC、TKY 和 CA 数据集上分别由 0.3505、0.2766 和 0.2060 提升至 0.3550、0.2829 和 0.2122，表现出较为稳定的改进趋势。进一步结合三组随机种子的重复实验结果可以看出，本文方法的平均性能增益具有一定稳定性，说明该改进并非偶然现象。
此外，本文还通过对比不同地理信息注入方式以及消融实验，进一步分析了各组成模块的作用。实验结果说明，相较于在生成阶段额外引入显式地理 token 的方式，将地理信息提前融入离散表示学习过程，更有利于提升 Semantic ID 对空间结构的表达能力，并最终带来更好的推荐性能。相关结果表明，地理信息不仅可以作为推荐模型推理阶段的辅助条件，更应当作为 POI 离散语义表示学习的重要组成部分加以建模。

本文的研究为生成式 POI 推荐中语义信息与空间信息的联合建模提供了新的思路，也为后续 Semantic ID 在位置推荐场景中的扩展应用奠定了一定基础。


\section{未来展望}
尽管本文的方法取得了一定效果，但仍有进一步完善的空间。
首先，本文当前对 POI 语义信息的建模主要依赖类别名称，语义来源相对有限，后续可以结合更丰富的文本描述、评论信息、区域功能属性以及多模态特征，提升 POI 表示的完整性。
其次，本文采用的地理增强方式主要面向静态空间结构，而用户真实的出行行为往往还受到时间周期、交通条件、活动目的和城市功能分区变化等动态因素影响，未来可以考虑将动态时空上下文进一步融入 Semantic ID 学习过程。与此同时，GeoPE 的参考点设置、RQ-VAE 的码本规模以及 SID 层次结构设计仍有优化空间，后续可从自适应空间划分、层级离散表示控制以及冲突率与推荐性能协同优化等角度继续改进。
除此之外，本文工作主要基于现有公开数据集完成验证，在跨城市迁移、冷启动 POI 建模以及更大规模真实场景部署方面仍需进一步实验支持。后续若能结合更丰富的用户行为数据、更强的生成式模型以及更稳定的空间语义联合表示机制，基于 Semantic ID 的生成式 POI 推荐方法仍具有较大的研究价值和应用潜力。